% 第 11 章 Working With Edges
\bichapter{处理边}{Working With Edges}
\label{chap:edges}

本章简要概述如何在 IC Validator runset 中使用边（edge）。边相关函数见 \textit{IC Validator Reference Manual}。

IC Validator 工具提供：
\begin{itemize}
  \item 将边层与父层关联
  \item 拓扑运算
  \item 逻辑运算
  \item 边操作
\end{itemize}

% ============================================================
\section{将边层与父层关联}
\subsection*{Associating Edge Layers With a Parent Layer}

边层可以有一个父多边形层（parent polygon layer）。
\begin{itemize}
  \item 所有边选择运算都会保持边层与其派生自的多边形层之间的连接。
  \item 所有边创建运算会得到“孤立”（orphan）边层，即没有父多边形层的边层。
  \item 父多边形层通过影响边的尺寸类函数行为，使输出更直观。
\end{itemize}
例如，图~\ref{fig:edge-parent} 展示边与父层的关联：
\begin{lstlisting}[style=pxl]
red = edge_derived_from_blue;
@ "red edge min width rule";
internal1(red, spacing < x);
\end{lstlisting}

IC Validator 工具会检测从 e2 到 e3 的违例，并避免从 e1 到 e3 的违例。

\figplaceholder{Figure 87: Edge Association With a Parent Layer}{边与父层的关联}{fig:edge-parent}

% ============================================================
\section{拓扑运算}
\subsection*{Topological Operations}

拓扑运算使用所选的完整边，且不添加新端点。这些运算包括：
\begin{itemize}
  \item \cmd{adjacent\_edge()} 与 \cmd{not\_adjacent\_edge()}
  \item \cmd{angle\_edge()} 与 \cmd{not\_angle\_edge()}
  \item \cmd{inside\_point\_touching\_edge()} 与 \cmd{not\_inside\_point\_touching\_edge()}
  \item \cmd{inside\_touching\_edge()} 与 \cmd{not\_inside\_touching\_edge()}
  \item \cmd{interacting\_edge()} 与 \cmd{not\_interacting\_edge()}
  \item \cmd{length\_edge()} 与 \cmd{not\_length\_edge()}
  \item \cmd{outside\_point\_touching\_edge()} 与 \cmd{not\_outside\_point\_touching\_edge()}
  \item \cmd{outside\_touching\_edge()} 与 \cmd{not\_outside\_touching\_edge()}
  \item \cmd{point\_touching\_edge()} 与 \cmd{not\_point\_touching\_edge()}
  \item \cmd{touching\_edge()} 与 \cmd{not\_touching\_edge()}
  \item \cmd{vertices\_edge()} 与 \cmd{not\_vertices\_edge()}
\end{itemize}
图~\ref{fig:touching-edges} 给出 touching edges 示例：
\begin{lstlisting}[style=pxl]
green = touching_edge (layer1 = blue, layer2 = red);
\end{lstlisting}

\figplaceholder{Figure 88: Touching Edges Example}{Touching edges 示例}{fig:touching-edges}

% ============================================================
\section{逻辑运算}
\subsection*{Logical Operations}

逻辑运算会生成新端点。这些运算包括：
\begin{itemize}
  \item \cmd{coincident\_edge()} 与 \cmd{not\_coincident\_edge()}
  \item \cmd{coincident\_inside\_edge()} 与 \cmd{not\_coincident\_inside\_edge()}
  \item \cmd{coincident\_outside\_edge()} 与 \cmd{not\_coincident\_outside\_edge()}
  \item \cmd{and\_edge()}
  \item \cmd{not\_edge()}
  \item \cmd{or\_edge()}
  \item \cmd{xor\_edge()}
\end{itemize}
图~\ref{fig:coincident-edges} 给出 coincident edges 示例：
\begin{lstlisting}[style=pxl]
green = coincident_edge (layer1 = blue, layer2 = red);
\end{lstlisting}

\figplaceholder{Figure 89: Coincident Edges Example}{Coincident edges 示例}{fig:coincident-edges}

% ============================================================
\section{边操作}
\subsection*{Edge Manipulation}

下列函数创建边；创建的边始终为孤立边（orphan edges）：
\begin{itemize}
  \item \cmd{move\_edge()}
  \item \cmd{snap\_edge()}
  \item \cmd{extend\_edge()}
  \item \cmd{copy\_edge()}
\end{itemize}
下列函数由边创建多边形：
\begin{itemize}
  \item \cmd{edge\_size()}
\end{itemize}
